\documentclass[9pt,a4paper]{extarticle}
\usepackage[T1]{fontenc}
\usepackage[utf8]{inputenc}
\usepackage{mathpazo}
\usepackage{courier}
\linespread{1.03}
\usepackage[a4paper,top=6mm,bottom=6mm,left=6mm,right=6mm]{geometry}
\setlength{\parindent}{0pt}\setlength{\parskip}{0pt}
\pagestyle{empty}

\usepackage[dvipsnames,svgnames,x11names]{xcolor}
\definecolor{clKw}  {HTML}{7a0000}
\definecolor{clTy}  {HTML}{0a3d6b}
\definecolor{clFn}  {HTML}{4a1a6b}
\definecolor{clSt}  {HTML}{1e4d2b}
\definecolor{clCm}  {HTML}{5a4e3a}
\definecolor{clOp}  {HTML}{1a0d00}
\definecolor{acBlue}  {HTML}{0a3d6b}
\definecolor{acTeal}  {HTML}{0a4e45}
\definecolor{acGreen} {HTML}{1e4d2b}
\definecolor{acPurple}{HTML}{4a1a6b}
\definecolor{acRed}   {HTML}{7a0000}
\definecolor{acAmber} {HTML}{7a3800}
\definecolor{acNavy}  {HTML}{0e1f4a}
\definecolor{acOrange}{HTML}{8a3200}
\definecolor{bgBlue}  {HTML}{f4f9fd}
\definecolor{bgTeal}  {HTML}{f0faf8}
\definecolor{bgGreen} {HTML}{f1f8f3}
\definecolor{bgPurple}{HTML}{f8f4fc}
\definecolor{bgRed}   {HTML}{fdf4f4}
\definecolor{bgAmber} {HTML}{fdf8ee}
\definecolor{bgNavy}  {HTML}{f3f5fb}
\definecolor{bgOrange}{HTML}{fdf4ee}
\definecolor{bgCode}  {HTML}{fefcfa}
\definecolor{bgHdr}   {HTML}{f7f3ee}
\definecolor{bdCard}  {HTML}{9a9080}
\definecolor{secBg}   {HTML}{1a0d00}

\usepackage{listings}
\lstdefinelanguage{CSharp}{
  sensitive=true,
  morekeywords=[1]{abstract,as,base,bool,break,byte,case,catch,char,checked,class,const,continue,decimal,default,delegate,do,double,else,enum,event,explicit,extension,extern,false,field,finally,fixed,float,for,foreach,goto,if,implicit,in,init,int,interface,internal,is,lock,long,namespace,new,not,null,object,operator,out,override,params,partial,private,protected,public,readonly,record,ref,return,sbyte,sealed,short,sizeof,stackalloc,static,string,struct,switch,this,throw,true,try,typeof,uint,ulong,unchecked,unsafe,ushort,using,virtual,void,volatile,when,where,while,with,nint,nuint,global,async,await,var,and,or,scoped,required,file,get,set,value,yield,allows},
  morekeywords=[2]{Action,Array,ArgumentException,ArrayPool,BinaryReader,BinaryWriter,BitConverter,BitOperations,CancellationToken,CollectionsMarshal,ConcurrentDictionary,Console,CultureInfo,Dictionary,Environment,Exception,File,FileStream,Func,GC,HashCode,HashSet,IComparable,IComparer,IDisposable,IEnumerable,IEnumerator,IEquatable,INumber,InvalidOperationException,IOException,IReadOnlyList,ISpanFormattable,ISpanParsable,Interlocked,KeyValuePair,List,Lock,Math,MathF,Memory,MemoryMarshal,NotImplementedException,Object,Parallel,ParallelOptions,Path,Queue,Random,ReadOnlyMemory,ReadOnlySpan,SearchValues,SortedDictionary,Span,Stack,Stopwatch,StreamReader,StreamWriter,String,StringBuilder,StringComparer,Task,Thread,ThreadLocal,Tuple,Type,Unsafe,ValueTask,ValueTuple,Vector,Vector128,Vector256},
  morekeywords=[3]{Abs,Add,AddRange,AllocateUninitializedArray,Any,AsSpan,AsMemory,AsRef,AsBytes,AsPointer,AsTask,Atan2,Cancel,Cast,Ceiling,Clear,Close,CompareExchange,Combine,Contains,CopyBlockUnaligned,CopyTo,Cos,CreateSpan,DangerousGetReference,Dispose,EnterScope,EnsureCapacity,Exists,Exp,Fill,Floor,For,ForEach,Format,GetHashCode,GetReference,GetValueRefOrAddDefault,Increment,IndexOf,IndexOfAny,IsBetween,Length,Log,Max,Min,Next,Open,OpenRead,OpenWrite,Parse,Peek,Pop,Pow,Push,ReadAllLines,ReadAllText,ReadLine,Rent,Remove,RemoveAt,Return,Reverse,Round,SequenceEqual,Sin,Slice,Sort,Sqrt,Sum,Tan,ThrowIfNull,ToArray,ToString,TryAdd,TryFormat,TryGetValue,TryParse,WaitAll,WhenAll,Write,WriteAllText,WriteAllLines,WriteLine,WriteUnaligned},
  morecomment=[l]{//},
  morecomment=[s]{/*}{*/},
  morestring=[b]",
  morestring=[b]',
}
\lstdefinestyle{cs}{
  language=CSharp,
  basicstyle=\ttfamily\bfseries\fontsize{6.25pt}{8.25pt}\selectfont,
  keywordstyle=[1]{\bfseries\color{clKw}},
  keywordstyle=[2]{\color{clTy}},
  keywordstyle=[3]{\color{clFn}},
  stringstyle=\color{clSt},
  commentstyle=\itshape\color{clCm},
  backgroundcolor=\color{bgCode},
  frame=none,
  xleftmargin=3pt, xrightmargin=2pt,
  aboveskip=1.5pt, belowskip=0.5pt,
  breaklines=true, keepspaces=true,
  showstringspaces=false, tabsize=2,
  literate=
    {=>}{{\textbf{\color{clOp}=>}}}{2}
    {->}{{\textbf{\color{clOp}->}}}{2}
    {??}{{\textbf{\color{clOp}??}}}{2}
    {?.}{{\textbf{\color{clOp}?.}}}{2}
    {>=}{{\textbf{\color{clOp}>=}}}{2}
    {<=}{{\textbf{\color{clOp}<=}}}{2}
    {==}{{\textbf{\color{clOp}==}}}{2}
    {!=}{{\textbf{\color{clOp}!=}}}{2}
    {=>}{{\textbf{\color{clOp}=>}}}{2},
}

\usepackage{tcolorbox}
\tcbuselibrary{skins,breakable}
\newenvironment{card}[3][acBlue]{%
  \begin{tcolorbox}[
    enhanced, colback=#2, colframe=bdCard,
    coltitle=black, colbacktitle=#2!70!bgHdr,
    fonttitle=\bfseries\fontsize{6.0pt}{7.2pt}\selectfont\MakeUppercase,
    title={#3},
    left=3pt, right=2pt, top=1.2pt, bottom=1.6pt,
    toptitle=1.0pt, bottomtitle=1.0pt,
    boxrule=0.4pt,
    borderline west={3pt}{0pt}{#1},
    arc=2pt, before skip=0pt, after skip=1.6mm,
    breakable,
  ]}{\end{tcolorbox}}

\usepackage{multicol}
\setlength{\columnsep}{2.8mm}
\newcommand{\sectionbanner}[1]{%
  \vspace{0.5mm}%
  {\setlength{\fboxsep}{1.5pt}%
   \colorbox{secBg}{\color{white}\bfseries\fontsize{6.5pt}{7.7pt}\selectfont
   \MakeUppercase{\hspace{0.5mm}#1\hspace{0.5mm}}}%
  }\vspace{1.3mm}\par\nobreak}
\newcommand{\sub}[1]{{\fontsize{5.5pt}{6.8pt}\selectfont\color{gray}\MakeUppercase{#1}}\par\vspace{1pt}}
\usepackage{booktabs,array,colortbl,enumitem}
\setlist[itemize]{leftmargin=3.8mm,itemsep=0.7pt,topsep=1pt,parsep=0pt,partopsep=0pt}
\newcommand{\code}[1]{{\ttfamily\bfseries #1}}

\begin{document}

{\centering
{\fontsize{13.0pt}{15.4pt}\selectfont\bfseries
  C\# for High-Performance Numerical and Scientific Code}\par\vspace{0.9mm}
{\fontsize{6.2pt}{7.3pt}\selectfont\color{gray}
  \MakeUppercase{types \textperiodcentered{} classes \textperiodcentered{} arrays \textperiodcentered{} spans \textperiodcentered{} refs \textperiodcentered{} memory layout \textperiodcentered{} hot loops \textperiodcentered{} deterministic APIs \textperiodcentered{} C\#~14 / .NET~10}}\par\vspace{1mm}
{\color{secBg}\rule{\linewidth}{1pt}}\par\vspace{1.8mm}}

%% ====================================================================
\sectionbanner{A Typical C\# Program}
\begin{multicols}{2}

\begin{card}[acNavy]{bgNavy}{Single-File Console Application}
\begin{lstlisting}[style=cs]
// File: Program.cs -- this is the entire program
using System.Globalization;

int n = 100;
double[] x = new double[n];
double[] y = new double[n];

for (int i = 0; i < n; i++)
  x[i] = Math.Sin(i * 0.1);

Scale(x, 2.0);
double d = Dot(x, x, n);
Console.WriteLine($"n = {n}, dot = {d:G10}");

// Static local functions:
static void Scale(Span<double> v, double a)
{
  for (int i = 0; i < v.Length; i++) v[i] *= a;
}

static double Dot(ReadOnlySpan<double> a,
    ReadOnlySpan<double> b, int n)
{
  double s = 0;
  for (int i = 0; i < n; i++) s += a[i] * b[i];
  return s;
}
\end{lstlisting}
\end{card}

\begin{card}[acTeal]{bgTeal}{What This Illustrates}
\begin{itemize}
\item \textbf{Top-level statements}: no \code{class}, no \code{Main}. The file \emph{is} the entry point.
\item \textbf{Static local functions}: defined after calling code, in the same file.
\item \textbf{Span parameters}: \code{Scale} takes a \code{Span<double>}---an array is implicitly convertible.
\item \textbf{String interpolation}: \code{\$"...\{d:G10\}"} for formatted output.
\end{itemize}
\sub{To compile and run}
\begin{lstlisting}[style=cs]
// $ dotnet new console -n MyApp
// $ cd MyApp && dotnet run -c Release
\end{lstlisting}
No headers, no makefiles, no linking step.
\end{card}

\end{multicols}

%% ====================================================================
\sectionbanner{Primitive Types, Math, and Syntax Essentials}
\begin{multicols}{3}

\begin{card}[acNavy]{bgNavy}{Numeric and Primitive Types}
\fontsize{6.0pt}{7.6pt}\selectfont
\begin{tabular}{@{}p{20mm}p{28mm}@{}}
\rowcolor{bgHdr}\textbf{Type} & \textbf{Description}\\
\code{int} / \code{long} & signed 32 / 64-bit integer\\
\code{float} / \code{double} & 32 / 64-bit IEEE~754\\
\code{decimal} & 128-bit exact decimal\\
\code{bool} & \code{true}/\code{false}\\
\code{byte}, \code{short} & unsigned 8 / signed 16-bit\\
\code{uint}, \code{ulong} & unsigned 32/64-bit\\
\code{nint}, \code{nuint} & platform-native size\\
\code{char} & UTF-16 character\\
\code{string} & immutable character sequence\\
\end{tabular}
\vspace{1pt}

All numeric types are fixed-size. Complex numbers via \code{System.Numerics.Complex}.
\end{card}

\begin{card}[acTeal]{bgTeal}{Control Flow Essentials}
\begin{lstlisting}[style=cs]
for (int i = 0; i < n; i++) { /* ... */ }
while (residual > tol) { /* ... */ }
do { /* ... */ } while (residual > tol);

if (x > 0) { /* ... */ }
else if (x == 0) { /* ... */ }
else { /* ... */ }

double y = (x > 0) ? x : -x; // ternary
\end{lstlisting}
\begin{itemize}
\item Statements end with \code{;}. Blocks use \code{\{~\}}. Case-sensitive.
\item \code{break} exits a loop; \code{continue} skips to next iteration.
\item \code{for (;;)} is an infinite loop; exit with \code{break} or \code{return}.
\end{itemize}
\end{card}

\begin{card}[acOrange]{bgOrange}{Math via \code{System.Math}}
\fontsize{6.0pt}{7.6pt}\selectfont
\begin{tabular}{@{}p{26mm}p{22mm}@{}}
\rowcolor{bgHdr}\textbf{Method} & \textbf{Description}\\
\code{Math.Sin}, \code{Cos}, \code{Tan} & trigonometric\\
\code{Math.Asin}, \code{Acos}, \code{Atan} & inverse trig\\
\code{Math.Atan2(y, x)} & two-argument arctangent\\
\code{Math.Sqrt(x)} & square root\\
\code{Math.Pow(a, b)} & $a^b$\\
\code{Math.Exp(x)}, \code{Log(x)} & $e^x$, $\ln x$\\
\code{Math.Abs(x)} & absolute value\\
\code{Math.Max(a,b)}, \code{Min} & extrema\\
\code{Math.Ceiling}, \code{Floor}, \code{Round} & rounding\\
\end{tabular}
\vspace{1pt}

\code{Math} operates on \code{double}; use \code{MathF} for \code{float}. No power operator; use \code{Math.Pow} or explicit multiplication.

Special: \code{double.NaN}, \code{double.PositiveInfinity}, \code{double.Epsilon}. Test with \code{double.IsNaN(x)}.
\end{card}

\end{multicols}

%% ====================================================================
\sectionbanner{C\# Language Foundations}
\begin{multicols}{3}

\begin{card}[acBlue]{bgBlue}{Namespaces, \code{using}, and Program Entry}
\begin{lstlisting}[style=cs]
// File: Program.cs  (top-level statements)
Console.WriteLine("Hello, World!");
\end{lstlisting}
Other files use file-scoped namespaces:
\begin{lstlisting}[style=cs]
namespace MyFEM;    // file-scoped (no braces)

using System.Numerics;          // import namespace
using static Math;              // import static members
using Vec = System.Numerics.Vector<double>; // alias
global using System.Globalization; // project-wide

public class Solver { /* ... */ }
\end{lstlisting}
\begin{itemize}
\item \code{using}: import a namespace into the current file.
\item \code{using static}: import static members directly (e.g.\ call \code{Sin(x)} instead of \code{Math.Sin(x)}).
\item \code{using Alias = Type}: create a short name for a long type.
\item \code{global using}: applies to every file in the project; place in one file.
\item Only one file may have top-level statements.
\end{itemize}
\end{card}

\begin{card}[acRed]{bgRed}{Access Modifiers}
\fontsize{6.0pt}{7.6pt}\selectfont
\begin{tabular}{@{}p{18mm}p{31mm}@{}}
\rowcolor{bgHdr}\textbf{Modifier} & \textbf{Visible to}\\
\code{public} & Everyone\\
\code{private} & Same class only (default for members)\\
\code{protected} & Same class + derived classes\\
\code{internal} & Same assembly (default for types)\\
\code{protected internal} & Assembly + derived classes\\
\code{private protected} & Same assembly + derived classes\\
\end{tabular}
\end{card}

\begin{card}[acGreen]{bgGreen}{Variables, Constants, \code{var}}
\begin{lstlisting}[style=cs]
int n = 100;
double tol = 1.0e-10;
float f = 3.14f;        // note f suffix
var x = 2.5;            // inferred as double
const double PI = 3.141592653589793;
const int MAX_ITER = 1000;
var list = new List<int>(); // type is List<int>
\end{lstlisting}
\begin{itemize}
\item \code{const}: compile-time. \code{readonly}: set once at construction.
\item \code{var} infers the type at compile time; it is not dynamic.
\end{itemize}
\end{card}

\begin{card}[acPurple]{bgPurple}{Classes: \code{field} Keyword and Properties}
\begin{lstlisting}[style=cs]
public class Material(double e, double nu)
{
  // Semi-auto property with 'field' (C# 14):
  public double E
  {
    get => field;
    set => field = value > 0
      ? value : throw new ArgumentException();
  } = e;

  public double Nu
  {
    get => field;
    set => field = value is > -1.0 and < 0.5
      ? value : throw new ArgumentException();
  } = nu;

  public double Density { get; set; } // auto
  public double Lambda { get; }       // read-only
    = e * nu / ((1 + nu) * (1 - 2 * nu));

  public double BulkModulus()
    => E / (3.0 * (1.0 - 2.0 * Nu));
}
\end{lstlisting}
\begin{itemize}
\item \code{field} (C\#~14) accesses the backing field in \code{get}/\code{set}, eliminating manual \code{\_field} declarations.
\item Expression-body \code{=>} for single-expression methods.
\end{itemize}
\end{card}

\begin{card}[acAmber]{bgAmber}{Primary Constructors and Initialization}
\begin{lstlisting}[style=cs]
public class Node(int id, double x, double y, double z)
{
  public int Id { get; } = id;
  public double X { get; set; } = x;
  public double Y { get; set; } = y;
  public double Z { get; set; } = z;
}
var n = new Node(5, 0, 0, 0) { X = 1.0, Y = 2.0 };
\end{lstlisting}
\begin{itemize}
\item Primary constructor parameters are in scope for initializers and method bodies.
\item Object initializers set properties after construction.
\end{itemize}
\end{card}

\begin{card}[acNavy]{bgNavy}{Inheritance and Virtual Methods}
\begin{lstlisting}[style=cs]
public abstract class Element(int id, int[] nodeIds)
{
  public int Id { get; } = id;
  public int[] NodeIds { get; } = nodeIds;
  public abstract double[,] Stiffness();
  public virtual int Order => 1;
}

public class Tri3(int id, int[] n) : Element(id, n)
{
  public override double[,] Stiffness()
    => throw new NotImplementedException();
  public override int Order => 1;
}
\end{lstlisting}
\begin{itemize}
\item \code{abstract}: must override; class cannot be instantiated.
\item \code{virtual}: can override. \code{base} calls the parent.
\item Single class inheritance; use interfaces for multiple contracts.
\end{itemize}
\end{card}

\end{multicols}

%% ====================================================================
\sectionbanner{Interfaces, Generics, Enums, and Pattern Matching}
\begin{multicols}{3}

\begin{card}[acTeal]{bgTeal}{Interfaces and Extension Members}
\begin{lstlisting}[style=cs]
public interface ISolver
{
  void Solve(ReadOnlySpan<double> rhs, Span<double> x);
  int Iterations { get; }
}

// A class can implement multiple interfaces:
public class PCG : ISolver, IDisposable
{
  public void Solve(ReadOnlySpan<double> rhs,
      Span<double> x) { /* ... */ }
  public int Iterations { get; private set; }
  public void Dispose() { /* cleanup */ }
}

// Extension members (C# 14):
extension(ReadOnlySpan<double> span)
{
  public double Norm2()
  {
    double s = 0;
    for (int i = 0; i < span.Length; i++)
      s += span[i] * span[i];
    return Math.Sqrt(s);
  }
}
\end{lstlisting}
\begin{itemize}
\item C\#~14 \code{extension} blocks replace the old \code{static class} pattern.
\item Interfaces can have default implementations since C\#~8.
\end{itemize}
\end{card}

\begin{card}[acOrange]{bgOrange}{Generics and \code{allows ref struct}}
\begin{lstlisting}[style=cs]
public class DenseVector<T>(int n)
    where T : INumber<T>
{
  private readonly T[] _data = new T[n];
  public int Length => _data.Length;
  public ref T this[int i] => ref _data[i];

  public T Norm()
  {
    T s = T.Zero;
    for (int i = 0; i < _data.Length; i++)
      s += _data[i] * _data[i];
    return T.CreateChecked(
      Math.Sqrt(double.CreateChecked(s)));
  }
}

// allows ref struct (C# 13):
static T Sum<T>(ReadOnlySpan<T> x)
    where T : INumber<T>, allows ref struct
{
  T s = T.Zero;
  for (int i = 0; i < x.Length; i++) s += x[i];
  return s;
}

static void Swap<T>(ref T a, ref T b)
  => (a, b) = (b, a);
\end{lstlisting}
\begin{itemize}
\item \code{allows ref struct} (C\#~13) lets generics accept \code{Span<T>} and other ref-like types.
\item JIT compiles specialized versions for each value type.
\end{itemize}
\end{card}

\begin{card}[acRed]{bgRed}{Enums}
\begin{lstlisting}[style=cs]
public enum ElementType
{ Tri3 = 3, Quad4 = 4, Tet4 = 5, Hex8 = 6 }

public enum BoundaryCondition
{ Free, Fixed, Prescribed }

[Flags]
public enum DOF
{
  None = 0, Ux = 1, Uy = 2, Uz = 4,
  Rx = 8, Ry = 16, Rz = 32,
  All = Ux | Uy | Uz | Rx | Ry | Rz,
}
DOF free = DOF.Ux | DOF.Uy;
bool hasUx = (free & DOF.Ux) != 0;
\end{lstlisting}
Enums are integers underneath. \code{[Flags]} enables bitwise combination.
\end{card}

\begin{card}[acBlue]{bgBlue}{Nullable Types and Null Safety}
\begin{lstlisting}[style=cs]
// Value types are non-nullable by default:
int n = 5;               // cannot be null
int? m = null;           // nullable value type

if (m.HasValue) Console.WriteLine(m.Value);
int k = m ?? 0;          // null-coalescing
string? name = null;
int len = name?.Length ?? 0; // null-conditional
x ??= new();             // assign only if null
\end{lstlisting}
\begin{itemize}
\item \code{T?} for value types wraps in \code{Nullable<T>}.
\item \code{??} / \code{?.} / \code{??=} for concise null handling.
\item Reference types are nullable by default; enable nullable context for warnings.
\end{itemize}
\end{card}

\begin{card}[acPurple]{bgPurple}{Pattern Matching and Switch Expressions}
\begin{lstlisting}[style=cs]
if (element is Tri3 tri)
  Console.WriteLine(tri.Order);

string name = elemType switch
{
  ElementType.Tri3  => "Triangle",
  ElementType.Quad4 => "Quadrilateral",
  ElementType.Tet4  => "Tetrahedron",
  _ => "Unknown",
};

string Classify(double x) => x switch
{
  < 0  => "negative",
  0    => "zero",
  < 1  => "small",
  >= 1 => "large",
};
\end{lstlisting}
The discard \code{\_} is the catch-all arm. Property patterns: \code{\{ E: > 200e9 \}}.
\end{card}

\begin{card}[acGreen]{bgGreen}{Tuples, Deconstruction, Records}
\begin{lstlisting}[style=cs]
var (det, singular) = ComputeDet(K);

static (double det, bool singular) ComputeDet(
    double[,] K)
{
  double d = /* ... */ 1.0;
  return (d, Math.Abs(d) < 1e-15);
}

// Record: value equality + ToString for free:
public record MaterialProps(double E, double Nu, double Rho);

// Record struct: stack-allocated:
public readonly record struct Vec2(double X, double Y);
\end{lstlisting}
\begin{itemize}
\item Tuples: lightweight multi-return without a dedicated type.
\item \code{record struct}: value semantics + equality + \code{ToString}.
\end{itemize}
\end{card}

\end{multicols}

%% ====================================================================
\sectionbanner{Arrays, Indexing, and Multidimensional Data}
\begin{multicols}{3}

\begin{card}[acAmber]{bgAmber}{One-Dimensional Arrays}
\begin{lstlisting}[style=cs]
double[] x = new double[100]; // zeroed
int[] ids = [0, 1, 2, 3];    // collection expression
x[0] = 1.0;
int n = x.Length;
Array.Resize(ref x, 200);    // new array
x.AsSpan(0, 100).CopyTo(y);
Array.Fill(x, 0.0);
Array.Sort(x);
\end{lstlisting}
\begin{itemize}
\item Arrays are heap-allocated reference types. 0-based indexing.
\item Bounds checked at runtime (\code{IndexOutOfRangeException}).
\item Collection expressions \code{[a, b]} work for arrays, lists, spans.
\end{itemize}
\end{card}

\begin{card}[acBlue]{bgBlue}{Multidimensional vs Jagged Arrays}
\begin{lstlisting}[style=cs]
double[,] K = new double[3, 3]; // row-major
K[0, 0] = 1.0;
int rows = K.GetLength(0);

double[][] J = new double[3][]; // jagged
J[0] = new double[3];

// Best for performance: flat 1D + manual index
double[] Kf = new double[n * n];
Kf[i * n + j] = value; // row-major
\end{lstlisting}
\begin{itemize}
\item \code{double[,]}: contiguous block. Jagged \code{[][]}: scattered rows.
\item Flat \code{double[]} with index arithmetic gives best cache locality.
\end{itemize}
\end{card}

\begin{card}[acTeal]{bgTeal}{Array Slicing and Ranges}
\begin{lstlisting}[style=cs]
double[] a = [1, 2, 3, 4, 5, 6, 7, 8];
double[] sub = a[2..5];  // ALLOCATES copy!
double last  = a[^1];    // 8 (index from end)

// Span slicing: no copy, no allocation
Span<double> view = a.AsSpan(2, 3);
ReadOnlySpan<double> ro = a.AsSpan()[2..5];
\end{lstlisting}
\begin{itemize}
\item \code{a[i..j]} on arrays allocates a new array. \code{AsSpan()[i..j]} does not.
\item \code{\^{}n} means ``n from the end''.
\item Always prefer \code{Span} slicing in hot paths.
\end{itemize}
\end{card}

\begin{card}[acRed]{bgRed}{\code{foreach}, LINQ, and When Not to Use Them}
\begin{lstlisting}[style=cs]
foreach (double xi in x) Console.WriteLine(xi);

// With index, use classic for:
for (int i = 0; i < x.Length; i++) x[i] *= 2.0;

// LINQ: declarative queries (ALLOCATES!)
double sum = x.Sum();
double max = x.Max();
double[] pos = x.Where(v => v > 0).ToArray();
int[] sorted = ids.OrderBy(i => coords[i]).ToArray();
\end{lstlisting}
\begin{itemize}
\item \code{foreach} cannot modify the collection.
\item LINQ allocates iterators, closures, and intermediate arrays.
\item \textbf{Never} use LINQ inside hot numerical loops.
\item LINQ is perfect for setup: mesh preprocessing, filtering, postprocessing.
\end{itemize}
\end{card}

\begin{card}[acNavy]{bgNavy}{Operator Overloading}
\begin{lstlisting}[style=cs]
public readonly struct Vec3(double x, double y, double z)
{
  public double X { get; } = x;
  public double Y { get; } = y;
  public double Z { get; } = z;

  public static Vec3 operator +(Vec3 a, Vec3 b)
    => new(a.X + b.X, a.Y + b.Y, a.Z + b.Z);
  public static Vec3 operator -(Vec3 a, Vec3 b)
    => new(a.X - b.X, a.Y - b.Y, a.Z - b.Z);
  public static Vec3 operator *(double s, Vec3 v)
    => new(s * v.X, s * v.Y, s * v.Z);
  public static Vec3 operator *(Vec3 v, double s)
    => s * v;
  public static double operator *(Vec3 a, Vec3 b)
    => a.X * b.X + a.Y * b.Y + a.Z * b.Z;

  public double Norm() => Math.Sqrt(this * this);
  public override string ToString()
    => $"({X:G6}, {Y:G6}, {Z:G6})";
}

// Usage reads like math:
Vec3 a = new(1, 0, 0);
Vec3 b = new(0, 1, 0);
Vec3 c = a + 2.0 * b;
double d = a * b; // dot product
\end{lstlisting}
Operators must be \code{static}. Both mixed-type orderings must be defined.
\end{card}

\begin{card}[acGreen]{bgGreen}{Indexers}
\begin{lstlisting}[style=cs]
public class DenseMatrix(int rows, int cols)
{
  private readonly double[] _d = new double[rows*cols];
  public int Rows => rows;
  public int Cols => cols;
  public ref double this[int i, int j]
    => ref _d[i * cols + j];
  public ReadOnlySpan<double> Row(int i)
    => _d.AsSpan(i * cols, cols);
}
var K = new DenseMatrix(3, 3);
K[0, 0] = 1.0;
\end{lstlisting}
Returning \code{ref double} allows in-place mutation without copying. This is the natural way to wrap flat storage as a 2D matrix. Use \code{Row()} to get zero-allocation views of individual rows.
\end{card}

\end{multicols}

%% ====================================================================
\sectionbanner{Exception Handling, Resources, and File I/O}
\begin{multicols}{3}

\begin{card}[acPurple]{bgPurple}{Exception Handling}
\begin{lstlisting}[style=cs]
try
{
  double[] K = LoadStiffness("mesh.dat");
  Solve(K, rhs, x);
}
catch (FileNotFoundException ex)
{
  Console.Error.WriteLine($"Not found: {ex.FileName}");
}
catch (ArithmeticException ex)
{
  Console.Error.WriteLine($"Math error: {ex.Message}");
}
catch (Exception ex)
{
  Console.Error.WriteLine($"Error: {ex.Message}");
  throw; // re-throw preserving stack trace
}
finally { /* always runs: cleanup code */ }
\end{lstlisting}
\begin{itemize}
\item Catch specific types first, general last.
\item \code{throw;} re-throws; \code{finally} always executes.
\item Exceptions in inner loops are expensive; validate beforehand.
\end{itemize}
\end{card}

\begin{card}[acAmber]{bgAmber}{\code{IDisposable} and \code{using}}
\begin{lstlisting}[style=cs]
using var reader = new StreamReader("mesh.dat");
while (reader.ReadLine() is { } line)
  ProcessLine(line);
// reader.Dispose() called at scope exit

// Custom disposable:
public class ScratchBuffer(int n) : IDisposable
{
  private double[]? _buf =
    ArrayPool<double>.Shared.Rent(n);
  public Span<double> Span => _buf.AsSpan();
  public void Dispose()
  {
    if (_buf is not null)
    { ArrayPool<double>.Shared.Return(_buf); _buf = null; }
  }
}
\end{lstlisting}
\begin{itemize}
\item \code{using var}: deterministic cleanup at scope exit.
\item Ideal for file handles, pooled buffers, unmanaged resources.
\end{itemize}
\end{card}

\begin{card}[acBlue]{bgBlue}{File I/O for Scientific Data}
\begin{lstlisting}[style=cs]
string[] lines = File.ReadAllLines("mesh.dat");
List<(int id, double x, double y, double z)> nodes = [];
foreach (string line in lines)
{
  ReadOnlySpan<char> s = line.AsSpan().Trim();
  if (s.IsEmpty || s[0] == '#') continue;
  var p = line.Split(' ',
    StringSplitOptions.RemoveEmptyEntries);
  nodes.Add((int.Parse(p[0]),
    double.Parse(p[1], CultureInfo.InvariantCulture),
    double.Parse(p[2], CultureInfo.InvariantCulture),
    double.Parse(p[3], CultureInfo.InvariantCulture)));
}

using var sw = new StreamWriter("results.dat");
sw.WriteLine("# Node  Ux  Uy");
for (int i = 0; i < n; i++)
  sw.WriteLine($"{i} {ux[i]:E14.7} {uy[i]:E14.7}");
\end{lstlisting}
\begin{itemize}
\item \code{StreamReader}/\code{StreamWriter} for large-file streaming.
\item Always use \code{CultureInfo.InvariantCulture} for numeric I/O.
\item Prefer \code{TryParse(ReadOnlySpan<char>)} to avoid string allocations.
\end{itemize}
\end{card}

\end{multicols}

%% ====================================================================
\sectionbanner{String Formatting, Delegates, and Console Output}
\begin{multicols}{3}

\begin{card}[acTeal]{bgTeal}{String Interpolation and Formatting}
\begin{lstlisting}[style=cs]
string s1 = $"x = {x:F6}";      // 3.141593
string s2 = $"x = {x:E4}";      // 3.1416E+000
string s3 = $"x = {x:G6}";      // 3.14159
string s4 = $"iter = {iter,6}";  // right-aligned
Console.WriteLine($"Step {iter}: res = {res:E4}");
\end{lstlisting}
\sub{Format specifiers}
\fontsize{6.0pt}{7.6pt}\selectfont
\begin{tabular}{@{}p{10mm}p{16mm}p{18mm}@{}}
\rowcolor{bgHdr}\textbf{Spec} & \textbf{Meaning} & \textbf{Example}\\
\code{F6} & fixed, 6 decimals & 3.141593\\
\code{E4} & scientific, 4 dec. & 3.1416E+000\\
\code{G6} & general, 6 digits & 3.14159\\
\code{N2} & thousands sep. & 1,234.56\\
\code{D5} & integer, padded & 00042\\
\end{tabular}
\end{card}

\begin{card}[acOrange]{bgOrange}{\code{StringBuilder} for Bulk Output}
\begin{lstlisting}[style=cs]
var sb = new StringBuilder(capacity: 1024);
sb.AppendLine("# Results");
for (int i = 0; i < n; i++)
{
  sb.Append(i).Append(' ');
  sb.Append(x[i].ToString("E10",
    CultureInfo.InvariantCulture));
  sb.AppendLine();
}
File.WriteAllText("output.dat", sb.ToString());
\end{lstlisting}
\begin{itemize}
\item \code{string} is immutable; loop concatenation is $O(n^2)$.
\item \code{StringBuilder} amortizes into a growable buffer.
\item For peak throughput, use \code{TryFormat} into span buffers.
\end{itemize}
\end{card}

\begin{card}[acNavy]{bgNavy}{Delegates, Lambdas, and \code{Func<>}}
\begin{lstlisting}[style=cs]
Func<double, double> f = Math.Sin;
Func<double, double> g = x => x * x - 2.0;
double root = Bisect(g, 0.0, 2.0, 1e-12);

static double Bisect(Func<double, double> f,
    double a, double b, double tol)
{
  while (b - a > tol)
  {
    double c = 0.5 * (a + b);
    if (f(a) * f(c) < 0) b = c; else a = c;
  }
  return 0.5 * (a + b);
}

Action<string> log = Console.WriteLine;
\end{lstlisting}
\begin{itemize}
\item \code{Func<..., TResult>}: returns a value. \code{Action<...>}: void.
\item Lambdas capture variables (not values); closures allocate.
\end{itemize}
\end{card}

\end{multicols}

%% ====================================================================
\sectionbanner{Mental Model and API Surface}
\begin{multicols}{3}

\begin{card}[acBlue]{bgBlue}{Performance Hierarchy}
\begin{itemize}
\item Fast code = fewer allocations, fewer cache misses, fewer indirect calls.
\item \code{for} in hot loops; LINQ for orchestration only.
\item Copies kill throughput with large structs or frequent slices.
\item Measure first: \code{Stopwatch} for coarse, BenchmarkDotNet for micro.
\end{itemize}
\sub{Rule of thumb}
Public API can be rich; internal kernels should be brutally explicit.
\end{card}

\begin{card}[acNavy]{bgNavy}{Choose Parameter and Return Types}
\fontsize{6.0pt}{7.6pt}\selectfont
\begin{tabular}{@{}p{19mm}p{31mm}@{}}
\rowcolor{bgHdr}\textbf{Type} & \textbf{Use when}\\
\code{T[]} & caller owns concrete array\\
\code{Span<T>} & mutable window, sync only\\
\code{ReadOnlySpan<T>} & read-only window, kernels\\
\code{Memory<T>} & buffer survives async boundary\\
\code{IReadOnlyList<T>} & indexed public API\\
\code{IEnumerable<T>} & lazy enumeration, not hot paths\\
\code{ref T} / \code{in T} & alias: mutable / read-only\\
\code{ref readonly T} & return alias without copy\\
\end{tabular}
\vspace{1pt}

Never expose \code{List<T>} unless caller mutation is intended. \code{params ReadOnlySpan<T>} (C\#~13) avoids implicit array allocation.
\end{card}

\begin{card}[acRed]{bgRed}{Struct vs Class}
\begin{itemize}
\item \code{struct}: value type. Assignment copies data. Stack or inline.
\item \code{class}: reference type. Assignment copies the pointer. Heap.
\item Good struct: tiny, copy-cheap, no identity, no inheritance.
\item Bad struct: large matrix, mutable graph node, owns arrays.
\item \code{readonly struct} for tiny immutable values. Pass large structs with \code{in}.
\end{itemize}
\begin{lstlisting}[style=cs]
public readonly struct Vec3(double x, double y, double z)
{
  public double X { get; } = x;
  public double Y { get; } = y;
  public double Z { get; } = z;
}

public readonly struct Mat6 // likely too big to copy
{
  private readonly double M00, M01, M02, ...;
}
\end{lstlisting}
\end{card}

\end{multicols}

%% ====================================================================
\sectionbanner{Spans, Memory Layout, and Pools}
\begin{multicols}{2}

\begin{card}[acGreen]{bgGreen}{Span, ReadOnlySpan, stackalloc, scoped}
\begin{lstlisting}[style=cs]
Span<double> tmp = stackalloc double[16];
tmp.Clear();

ReadOnlySpan<char> token = line.AsSpan(start, len);
double.TryParse(token, System.Globalization.NumberStyles.Float,
    CultureInfo.InvariantCulture, out double x);

static int CountPositives(scoped ReadOnlySpan<double> x)
{
  int c = 0;
  for (int i = 0; i < x.Length; ++i)
    if (x[i] > 0.0) ++c;
  return c;
}

// params span (C# 13): no hidden array allocation
static double Max(params ReadOnlySpan<double> values)
{
  double m = double.NegativeInfinity;
  foreach (double v in values) if (v > m) m = v;
  return m;
}
\end{lstlisting}
\begin{itemize}
\item \code{Span<T>}: stack-only, cannot be boxed or cross \code{await}. Think pointer~+ length.
\item \code{stackalloc}: tiny scratch buffers. \code{scoped}: prevents accidental escape.
\item Prefer \code{a.AsSpan(2, 6)} over \code{a[2..8]} in hot paths (no copy).
\end{itemize}
\end{card}

\begin{card}[acPurple]{bgPurple}{MemoryMarshal and Unsafe}
\begin{lstlisting}[style=cs]
Span<byte> bytes = stackalloc byte[64];
Span<double> x = MemoryMarshal.Cast<byte, double>(bytes);
ref double r0 = ref MemoryMarshal.GetReference(x);

unsafe
{
  fixed (double* p = x)
    for (int i = 0; i < x.Length; ++i) p[i] = 0.0;
}
\end{lstlisting}
\begin{itemize}
\item \code{MemoryMarshal.Cast<A,B>}: only when layout is guaranteed.
\item \code{CollectionsMarshal.AsSpan(list)}: valid only while list is not resized.
\item \code{Unsafe.Add(ref r0, i)} can remove bounds checks, but only when invariants are airtight.
\item Prefer explicit layout reasoning over ``clever'' reinterpretation.
\item \code{unsafe} requires \code{<AllowUnsafeBlocks>true</AllowUnsafeBlocks>}.
\end{itemize}
\end{card}

\begin{card}[acOrange]{bgOrange}{Pools and Scratch Space}
\begin{lstlisting}[style=cs]
var pool = ArrayPool<double>.Shared;
double[] tmp = pool.Rent(n);
try
{
  var span = tmp.AsSpan(0, n);
  span.Clear();
  Work(span);
}
finally
{
  pool.Return(tmp, clearArray: false);
}
\end{lstlisting}
\begin{itemize}
\item Rent when buffers are too large for \code{stackalloc}.
\item Never assume rented arrays are zeroed. Return exactly once.
\item For long-lived ownership, allocate your own array instead.
\item Think of pooled arrays as pre-allocated work space that avoids GC pressure.
\end{itemize}
\end{card}

\begin{card}[acTeal]{bgTeal}{Low-Level Memory}
\begin{itemize}
\item Flat arrays beat \code{T[,]} and jagged \code{T[][]} in kernels.
\item Zero-initialization is not free; skip it when all values will be overwritten.
\item Pinning blocks GC movement; keep pin scopes short.
\item Arrays $> 85$\,kB land on the Large Object Heap; prefer \code{ArrayPool}.
\item Locality dominates: contiguous traversal beats random access.
\item The GC is generational; short-lived small objects are cheap.
\end{itemize}
\begin{lstlisting}[style=cs]
static int Idx(int i, int j, int ncols) => i * ncols + j;

// Uninitialized (all values will be overwritten):
double[] buf = GC.AllocateUninitializedArray<double>(n);

// Pinned (won't move during GC):
double[] pin = GC.AllocateArray<double>(n, pinned: true);
\end{lstlisting}
\end{card}

\end{multicols}

%% ====================================================================
\sectionbanner{Collections, APIs, and Hot-Loop Discipline}
\begin{multicols}{3}

\begin{card}[acAmber]{bgAmber}{Collections in Kernels}
\begin{lstlisting}[style=cs]
var list = new List<int>(capacity: n);
list.Add(1);
list.EnsureCapacity(2 * n);

Span<int> s = CollectionsMarshal.AsSpan(list);
for (int i = 0; i < s.Length; ++i) s[i] *= 2;
\end{lstlisting}
\begin{itemize}
\item Pre-size \code{List<T>} and \code{Dictionary<K,V>} when cardinality is known.
\item \code{TryGetValue} beats contains-then-index. \code{HashSet<T>} for membership.
\item Do not mutate a list after taking its marshal span.
\end{itemize}
\sub{Essential collections}
\fontsize{6.0pt}{7.6pt}\selectfont
\begin{tabular}{@{}p{17mm}p{32mm}@{}}
\rowcolor{bgHdr}\textbf{Type} & \textbf{Use case}\\
\code{List<T>} & Growable array\\
\code{Dictionary<K,V>} & Hash map (sparse lookups)\\
\code{HashSet<T>} & Unique membership\\
\code{Queue<T>} & FIFO buffer (BFS, front advance)\\
\code{Stack<T>} & LIFO buffer (DFS, undo)\\
\code{SortedDictionary} & Ordered keys\\
\end{tabular}
\end{card}

\begin{card}[acBlue]{bgBlue}{Dictionary by Reference}
\begin{lstlisting}[style=cs]
ref double value = ref CollectionsMarshal
  .GetValueRefOrAddDefault(map, key, out bool exists);
if (!exists) value = 0.0;
value += increment;
\end{lstlisting}
\begin{itemize}
\item Avoids repeated hashing and extra lookups.
\item The returned ref is invalid if the dictionary resizes.
\item Excellent for sparse assembly-like accumulation.
\end{itemize}
\end{card}

\begin{card}[acNavy]{bgNavy}{Span-Based Parsing}
\begin{lstlisting}[style=cs]
ReadOnlySpan<char> line = source.AsSpan();
int k = line.IndexOf(';');
ReadOnlySpan<char> left  = line[..k];
ReadOnlySpan<char> right = line[(k + 1)..];
int id = int.Parse(left, CultureInfo.InvariantCulture);
double x = double.Parse(right, CultureInfo.InvariantCulture);
\end{lstlisting}
\begin{itemize}
\item Always \code{InvariantCulture} for persisted numeric data.
\item \code{SearchValues<char>} and \code{IndexOfAny} for high-volume tokenization.
\end{itemize}
\end{card}

\begin{card}[acGreen]{bgGreen}{Equality, Hashing, and Stable Keys}
\begin{itemize}
\item If a value is used as a key, its equality and hash code must be stable for its entire dictionary residency.
\item Override \code{Equals} and \code{GetHashCode} together.
\item Use ordinal string comparison for identifiers: \code{StringComparer.Ordinal}.
\item For floating-point keys, think hard before using raw doubles.
\end{itemize}
\begin{lstlisting}[style=cs]
public readonly struct Edge(int a, int b)
    : IEquatable<Edge>
{
  public int A { get; } = a;
  public int B { get; } = b;
  public bool Equals(Edge o) => A == o.A && B == o.B;
  public override bool Equals(object? obj)
    => obj is Edge e && Equals(e);
  public override int GetHashCode()
    => HashCode.Combine(A, B);
  public static bool operator ==(Edge l, Edge r)
    => l.Equals(r);
  public static bool operator !=(Edge l, Edge r)
    => !l.Equals(r);
}
\end{lstlisting}
Full contract: \code{IEquatable<T>}, \code{Equals(object?)}, \code{GetHashCode()}, \code{==}/\code{!=}. Or just use \code{record struct} to get all for free.
\end{card}

\begin{card}[acPurple]{bgPurple}{Public API Layering}
\fontsize{6.0pt}{7.4pt}\selectfont
\begin{tabular}{@{}p{18mm}p{31mm}@{}}
\rowcolor{bgHdr}\textbf{Layer} & \textbf{Typical signature}\\
Convenience & \code{double[] Solve(double[] rhs)}\\
Core sync & \code{void Solve(ReadOnlySpan<double> rhs, Span<double> x)}\\
Batched & \code{void SolveBatch(ReadOnlySpan<double> A, ...)}\\
Async I/O & \code{ValueTask<int> ReadAsync(Memory<byte> dst)}\\
\end{tabular}
\vspace{2pt}
Expose friendly wrappers, but keep one explicit zero-allocation kernel underneath.
\end{card}

\begin{card}[acRed]{bgRed}{Closure Trap}
\begin{lstlisting}[style=cs]
// BUG: all tasks capture the *same* variable i
for (int i = 0; i < n; ++i)
  tasks[i] = Task.Run(() => Work(i));

// FIX: copy to a local before capturing
for (int i = 0; i < n; ++i)
{
  int j = i;
  tasks[i] = Task.Run(() => Work(j));
}
\end{lstlisting}
\begin{itemize}
\item Lambdas capture \emph{variables}, not values.
\item Closures also cause hidden heap allocations for captured state.
\end{itemize}
\end{card}

\end{multicols}

%% ====================================================================
\sectionbanner{SIMD, Parallelism, and Diagnostics}
\begin{multicols}{2}

\begin{card}[acOrange]{bgOrange}{SIMD and Intrinsics}
\begin{lstlisting}[style=cs]
if (Vector.IsHardwareAccelerated)
{
  int w = Vector<double>.Count;
  int i = 0;
  Vector<double> acc = Vector<double>.Zero;
  for (; i <= n - w; i += w)
    acc += new Vector<double>(a, i) * new Vector<double>(b, i);
  double s = 0.0;
  for (int j = 0; j < w; ++j) s += acc[j];
  for (; i < n; ++i) s += a[i] * b[i];
}
\end{lstlisting}
\begin{itemize}
\item Start with \code{Vector<T>} before AVX intrinsics.
\item \code{Vector128<T>} / \code{Vector256<T>} for explicit width control.
\item Data layout often matters more than instruction cleverness.
\end{itemize}
\end{card}

\begin{card}[acNavy]{bgNavy}{Parallel Loops}
\begin{lstlisting}[style=cs]
Parallel.For(0, n, i => y[i] = F(x[i]));

Lock sync = new(); // System.Threading.Lock (C# 13)
double total = 0.0;
Parallel.For<double>(0, n,
  () => 0.0,
  (i, _, local) => local + a[i],
  local => { lock (sync) total += local; });

var opts = new ParallelOptions
{ MaxDegreeOfParallelism = Environment.ProcessorCount };
Parallel.For(0, nElem, opts, e =>
{
  Span<double> Ke = stackalloc double[ndof * ndof];
  ComputeElement(e, Ke);
  lock (sync) AssembleGlobal(e, Ke, K);
});
\end{lstlisting}
\begin{itemize}
\item \code{Lock} (C\#~13) replaces \code{object}-as-lock with a purpose-built, more efficient type.
\item Parallelize only when per-iteration work is large enough.
\item Avoid shared writes; partition by output slices.
\item Beware false sharing when adjacent elements are written by different threads.
\item Deterministic scientific code usually wants explicit partitioning.
\end{itemize}
\end{card}

\begin{card}[acBlue]{bgBlue}{Diagnostics and Measurement}
\begin{lstlisting}[style=cs]
Span<char> dst = stackalloc char[64];
if (value.TryFormat(dst, out int written,
    provider: CultureInfo.InvariantCulture))
  Console.WriteLine(dst[..written].ToString());

var sw = Stopwatch.StartNew();
for (int i = 0; i < nReps; i++) Solve(K, rhs, x);
sw.Stop();
Console.WriteLine(
  $"Avg: {sw.Elapsed.TotalMilliseconds / nReps:F2} ms");
\end{lstlisting}
\begin{itemize}
\item Warm up before timing. Release build only. No debugger.
\item Track allocations as a first-class metric, not only elapsed time.
\item Tools: \code{dotnet-counters}, \code{dotnet-trace}, BenchmarkDotNet.
\end{itemize}
\sub{Benchmark checklist}
\begin{itemize}
\item Stable input, pinned CPU frequency if possible, multiple iterations, verify outputs.
\end{itemize}
\end{card}

\end{multicols}

\end{document}
